\section{Dynamic Betti numbers are equivalent to dynamic rank}\label{sec:equiv}

This section proves that dynamic Betti numbers and dynamic rank of matrices with entries in $\{0,\pm1\}$ reduce to each other with constant overhead and a linear change in size (Theorem~\ref{thm:equiv} and Corollary~\ref{cor:equiv}); over $\Ftwo$ and over $\mathbb{F}_3$, every matrix has entries in $\{0,\pm1\}$.
Section~\ref{sec:b2r} gives the easy direction, from Betti numbers to rank.
The other direction rests on the gadget of Section~\ref{sec:gadget}.
Every row of the matrix becomes a hollow triangle, every column becomes a punctured sphere assembled from constant-size pieces, and each entry decides whether a hole of the sphere of its column is closed off or joined by a tube to the triangle of its row (Figure~\ref{fig:gadget}).
Section~\ref{sec:crux} proves that the $2$-cycles of the resulting complex correspond exactly to the linear dependencies among the active columns, restricted to the active rows, and Section~\ref{sec:equiv-thm} assembles the reductions.

\subsection{From Betti numbers to rank}\label{sec:b2r}

Let $\cU$ be a universe and $p\ge0$.
For $q\ge1$ let $B_q$ be the boundary matrix of $\cU$ in dimension~$q$, with rows indexed by $\cU_{q-1}$ and columns by~$\cU_q$.
For a subcomplex $K\subseteq\cU$ let $B_q^K$ be the matrix obtained from $B_q$ by replacing the columns of the $q$-simplices not in $K$ with zero columns.

\begin{lemma}\label{lem:betti-to-rank}
For every subcomplex $K\subseteq\cU$ and every $p\ge0$,
\[
  \beta_p(K)=n_p(K)-\rk\begin{pmatrix}B_p^K&0\\0&B_{p+1}^K\end{pmatrix},
\]
where $B_0^K$ is the empty matrix.
Inserting or deleting a $q$-simplex with $q\ge1$ changes one column of $B_q^K$, and no other column of $B_p^K$ or $B_{p+1}^K$; inserting or deleting a vertex changes only~$n_0(K)$.
\end{lemma}

\begin{proof}
The columns of $B_q^K$ indexed by $K_q$ form the matrix of $\bd_q$ on $K$, extended by rows for the $(q-1)$-simplices of $\cU$ that are not in $K$.
These extra rows are zero in those columns, because every facet of a simplex of $K$ lies in~$K$.
Adding zero rows and zero columns does not change the rank, so $\rk B_q^K=\rk\bd_q$, and the formula follows from~\eqref{eq:betti-rank}.
The column of a $q$-simplex $\sigma$ is zero in $B_q^K$ if $\sigma\notin K$ and equals its column in $B_q$ otherwise, so updating $\sigma$ changes exactly that column.
\end{proof}

A column change of $B_q^K$ is a single switch, or at most $q+1$ entry updates.
We call \emph{dynamic rank under switches} the variant of Problem~\ref{prob:rank} whose instance also contains a matrix $A_0\in\F^{r\times k}$, with $P=\supp(A_0)$, and in which $A=A_0$ throughout, so that the only updates are switches; its size is again $N=r+k+m$.
So for fixed $p$, dynamic $\beta_p$ reduces to dynamic rank under switches with overhead $(1,1)$ and to dynamic rank with overhead $(p+2,1)$: one data structure for the block matrix of Lemma~\ref{lem:betti-to-rank} and a counter for $n_p(K)$ suffice.
Combined with the fastest known dynamic-rank algorithm, this gives a randomized data structure with worst-case update time $O(n^{1.407})$ for every Betti number (Corollary~\ref{cor:upper} in Section~\ref{sec:lower-worst}).

\subsection{The gadget}\label{sec:gadget}

We now go the other way.
Fix a $0/1$ matrix $A$ with $r$ rows and $k$ columns.
We build a $2$-complex whose $2$-cycles encode the kernel of~$A$.
The construction uses four building blocks, each attached to \emph{triangle cycles}.
A triangle cycle $a=(a_0,a_1,a_2)$ consists of three vertices and the three edges $a_0a_1$, $a_1a_2$, $a_2a_0$; we call the triangle $a_0a_1a_2$ its \emph{filling}.
Indices of vertices of a triangle cycle are taken modulo~$3$, and we write $uv$ for the edge $\{u,v\}$ and $uvw$ for the triangle $\{u,v,w\}$.
\begin{itemize}
\item The \emph{cone} on $a$ with apex $o$ is $\{o\,a_t\,a_{t+1} : t\in\Z_3\}$, a disk with boundary~$a$ (three triangles).
\item The \emph{cylinder} between disjoint triangle cycles $a$ and $b$ is $\mathrm{Cyl}(a,b)=\{a_t\,a_{t+1}\,b_t,\; a_{t+1}\,b_t\,b_{t+1} : t\in\Z_3\}$, an annulus with boundary $a\sqcup b$ (six triangles).
\item The \emph{punctured cylinder} $\mathrm{Cyl}^{\circ}(a,b)=\mathrm{Cyl}(a,b)\setminus\{h\}$ is the cylinder with the triangle \mbox{$h=a_0a_1b_0$} removed (five triangles). We call the triangle cycle $(a_0,a_1,b_0)$ its \emph{hole}.
\item The \emph{tube} from the hole $h$ to a triangle cycle $c$ is $\mathrm{Cyl}(h,c)$ (six triangles).
\end{itemize}
A \emph{socket} between $a$ and $b$ consists of the punctured cylinder $\mathrm{Cyl}^{\circ}(a,b)$ and one or two \emph{attachments} at its hole: the \emph{seal}, which is the triangle~$h$, and the \emph{tube} $\mathrm{Cyl}(h,c)$ to a triangle cycle~$c$ disjoint from $a$ and~$b$.
In every complex $K(A;R,J)$ below, exactly one attachment of each socket is present.
With the seal, the socket is the annulus $\mathrm{Cyl}(a,b)$, and we call it \emph{sealed}; with the tube, it is a \emph{pair of pants}, a sphere with three holes, with boundary $a\sqcup b\sqcup c$ (see the proof of Lemma~\ref{lem:surface}), and we call it \emph{tubed}.
Figure~\ref{fig:gadget} shows the whole construction schematically, and Figure~\ref{fig:socket} shows a socket and one column surface.


% Figure: the pieces of the gadget (Definition~\ref{def:gadget}): a socket and one column surface.
% Requires tikz with the patterns library. Okabe--Ito colors: rows vermillion, columns blue.
\begin{figure}[t]
\centering
\definecolor{gRow}{HTML}{D55E00}%
\definecolor{gCol}{HTML}{0072B2}%
\colorlet{gLine}{gCol!80!black}%
\tikzset{
  gsurf/.style={fill=gCol!10,draw=gLine,line width=0.6pt},
  gtube/.style={double=gCol!24,double distance=4.4pt,line width=0.55pt,draw=gLine},
  growcyc/.style={draw=gRow,line width=1.15pt,line join=round},
  gseal/.style={fill=black!72,draw=black,line width=0.4pt},
  ghole/.style={fill=white,draw=gLine,line width=0.5pt},
  gteth/.style={fill=black!28,draw=black!80,line width=0.75pt,line join=round},
  glab/.style={font=\small},
  gsmall/.style={font=\footnotesize},
  gpanel/.style={font=\small}}
\providecommand\Description[2][]{}%
\Description{Left: a cylinder between two triangle cycles, cut open into a strip of six triangles, one of which is removed and hatched; below it the two ways to close the hole, a single dark triangle (seal) or a strip of six triangles (tube). Right: one column surface drawn as a long capsule divided by six triangle cycles into a cone, a socket tubed downward to the tether, four sockets of which two are tubed up to row cycles and two are sealed, and a cone.}
\begin{tikzpicture}[x=1cm,y=1cm,baseline=(pa.base)]
  %%% (b) one socket, cut open along the edge a_0 b_0
  \begin{scope}[shift={(0,0.95)}]
    \def\dx{0.82}\def\hy{1.2}
    \foreach \t in {0,1,2,3} { \coordinate (a\t) at (\t*\dx,\hy); \coordinate (b\t) at (\t*\dx,0); }
    \fill[gCol!10] (a0) -- (a3) -- (b3) -- (b0) -- cycle;
    \fill[white] (a0) -- (a1) -- (b0) -- cycle;
    \fill[pattern=north east lines,pattern color=black!45] (a0) -- (a1) -- (b0) -- cycle;
    \foreach \t/\u in {0/1,1/2,2/3} \draw[gLine,line width=0.5pt] (a\u) -- (b\t);
    \foreach \t in {1,2} \draw[gLine,line width=0.5pt] (a\t) -- (b\t);
    \draw[gLine,line width=1pt] (a0) -- (a3); \draw[gLine,line width=1pt] (b0) -- (b3);
    \draw[gLine,densely dashed,line width=0.6pt] (a0) -- (b0); \draw[gLine,densely dashed,line width=0.6pt] (a3) -- (b3);
    \foreach \t/\l in {0/0,1/1,2/2,3/0} { \node[gsmall,above] at (a\t) {$a_{\l}$}; \node[gsmall,below] at (b\t) {$b_{\l}$}; }
    \node[gsmall,fill=white,inner sep=1pt] at (0.28,0.8) {$h$};
    % the two attachments of the hole h
    \begin{scope}[shift={(0.05,-1.25)}]
      \draw[gseal] (0.0,0.0) -- (0.55,0.0) -- (0.0,0.45) -- cycle;
      \node[gsmall,align=center] at (0.3,-0.42) {seal $h$\\[-1.5pt]\textcolor{black!60}{1 triangle}};
      \node[gsmall] at (1.03,0.22) {or};
      \begin{scope}[shift={(1.45,0)}]
        \foreach \t in {0,1,2,3} { \coordinate (h\t) at (\t*0.3,0.45); \coordinate (c\t) at (\t*0.3,0); }
        \fill[gCol!24] (h0) -- (h3) -- (c3) -- (c0) -- cycle;
        \foreach \t/\u in {0/1,1/2,2/3} \draw[gLine,line width=0.45pt] (h\u) -- (c\t);
        \foreach \t in {0,1,2,3} \draw[gLine,line width=0.45pt] (h\t) -- (c\t);
        \draw[gLine,line width=0.8pt] (h0) -- (h3);
        \draw[gRow,line width=1pt] (c0) -- (c3);
        \node[gsmall,align=center] at (0.45,-0.42) {tube to $C_i$\\[-1.5pt]\textcolor{black!60}{6 triangles}};
      \end{scope}
    \end{scope}
    \node[gpanel] (pa) at (1.23,-2.45) {(a)};
  \end{scope}
\end{tikzpicture}\hfill
\begin{tikzpicture}[x=0.8cm,y=0.85cm,baseline=(pb.base)]
  %%% (c) one column surface: cone, tether socket, one socket per admissible position, cone
  \draw[gtube] (1.6,0.42) -- (1.6,-0.98);
  \draw[gtube] (3.6,0.42) -- (3.6,1.32);
  \draw[gtube] (9.6,0.42) -- (9.6,1.32);
  \draw[gsurf] (0.6,0.5) -- (10.6,0.5) arc (90:-90:0.62 and 0.5) -- (0.6,-0.5) arc (270:90:0.62 and 0.5) -- cycle;
  \foreach \x/\t in {0.6/0,2.6/1,4.6/2,6.6/3,8.6/4,10.6/5} {
    \draw[gLine,densely dashed,line width=0.45pt] (\x,0.5) arc (90:-90:0.17 and 0.5);
    \draw[gLine,line width=0.45pt] (\x,0.5) arc (90:270:0.17 and 0.5);
    \node[gsmall] at (\x,-0.76) {$s^j_{\t}$};
  }
  \draw[ghole] (1.6,-0.47) ellipse (0.14 and 0.05);
  \draw[ghole] (3.6,0.47) ellipse (0.14 and 0.05);
  \draw[gseal] (5.6,0.47) ellipse (0.15 and 0.06);
  \draw[gseal] (7.6,0.47) ellipse (0.15 and 0.06);
  \draw[ghole] (9.6,0.47) ellipse (0.14 and 0.05);
  \draw[growcyc] (3.24,1.32) -- (3.96,1.32) -- (3.6,1.84) -- cycle; \node[gsmall,text=gRow,anchor=west] at (4.0,1.62) {$C_{i_1}$};
  \draw[growcyc] (9.24,1.32) -- (9.96,1.32) -- (9.6,1.84) -- cycle; \node[gsmall,text=gRow,anchor=west] at (10.0,1.62) {$C_{i_4}$};
  \draw[gteth] (1.3,-0.98) -- (1.9,-0.98) -- (1.6,-1.44) -- cycle; \node[gsmall,anchor=west] at (1.98,-1.2) {$\Lambda_j$};
  \node[gsmall,text=black!65,anchor=east] at (-0.08,0.0) {cone};
  \node[gsmall,text=black!65,anchor=west] at (11.3,0.0) {cone};
  \foreach \xm/\t in {5.6/{$A_{i_2j}=0$},7.6/{$A_{i_3j}=0$}} \node[gsmall,text=black!65] at (\xm,0.95) {\t};
  \node[gpanel] (pb) at (5.6,-1.6) {(b)};
\end{tikzpicture}
\caption{The pieces of the gadget (Definition~\ref{def:gadget}).
(a)~A socket: the punctured cylinder, that is, the cylinder $\mathrm{Cyl}(a,b)$ between the triangle cycles $a$ and~$b$, cut open along the edge $a_0b_0$ (dashed; the two copies are identified), with the triangle $h=a_0a_1b_0$ removed (hatched), and its two possible attachments.
The hole~$h$ is either sealed by the single triangle~$h$ or tubed to a row cycle~$C_i$ by the six triangles of $\mathrm{Cyl}(h,C_i)$, so flipping one entry of~$A$ costs $1+6=7$ triangle updates (Lemma~\ref{lem:locality}).
(b)~The surface $\Sigma_j$ formed by the triangles of column~$j$: from left to right, a cone, the socket tubed to the tether~$\Lambda_j$, one socket per admissible position of column~$j$ (here four, in rows $i_1<i_2<i_3<i_4$, with $A_{i_1j}=A_{i_4j}=1$ and $A_{i_2j}=A_{i_3j}=0$), and a cone, glued in a chain along the triangle cycles $s^j_0,\dots,s^j_5$.}
\label{fig:socket}
\end{figure}


\begin{definition}[The gadget]\label{def:gadget}
Let $P\subseteq[r]\times[k]$ be a set of admissible positions, and for $j\in[k]$ let $S_j=\{i:(i,j)\in P\}$, listed in increasing order as $i_1<\dots<i_{d_j}$.
The complex $\cU(P)$ consists of the following simplices, all on distinct vertices unless stated otherwise.
\begin{itemize}
\item For each row $i\in[r]$, a triangle cycle $C_i$, the \emph{row cycle}, and its filling $D_i$.
\item For each column $j\in[k]$, a triangle cycle $\Lambda_j$, the \emph{tether}, and its filling $E_j$.
\item For each column $j$, triangle cycles $s^j_0,s^j_1,\dots,s^j_{d_j+1}$, cones on $s^j_0$ and on $s^j_{d_j+1}$, a socket between $s^j_0$ and $s^j_1$ whose only attachment is the tube to~$\Lambda_j$, and, for every $t\in[d_j]$, a socket between $s^j_t$ and $s^j_{t+1}$ with both attachments, the seal and the tube to $C_{i_t}$, which we call the \emph{socket of position} $(i_t,j)$.
\item All vertices and edges of the simplices above.
\end{itemize}
For a $0/1$ matrix $A$ with $\supp(A)\subseteq P$ and sets $R\subseteq[r]$, $J\subseteq[k]$, the complex $K(A;R,J)\subseteq\cU(P)$ consists of all vertices and edges of $\cU(P)$, the cones, the punctured cylinders, the tubes to the tethers, and
\begin{itemize}
\item for every $(i,j)\in P$, the seal of the socket of position $(i,j)$ if $A_{ij}=0$ and its tube to $C_i$ if $A_{ij}=1$;
\item the filling $D_i$ for every row $i\notin R$, and the filling $E_j$ for every column $j\in J$.
\end{itemize}
\end{definition}

So row $i$ is \emph{active} when its row cycle is hollow, that is, when its filling $D_i$ is absent, and column $j$ is \emph{active} when its tether is filled.
We write $\KA=K(A;[r],[k])$ for the complex in which all rows and all columns are active.
Row cycles and punctured spheres already appear in the static reduction~\cite[Theorem~3.1]{EP14}; the constant-size sockets and the switches make every matrix operation local (Lemma~\ref{lem:locality}).
We write $\Sigma_j$ for the set of triangles of $K(A;R,J)$ that belong to column~$j$: its two cones, its punctured cylinders, its seals and tubes.
Only admissible positions get sockets; a socket at a position that is always zero would always be sealed and would contribute nothing.
The size of $\cU(P)$ is $O(r+k+|P|)$, and $\cU(P)$ can be constructed in linear time.

\begin{lemma}\label{lem:surface}
Let $A$ be a $0/1$ matrix with $\supp(A)\subseteq P$, let $R\subseteq[r]$ and $J\subseteq[k]$, and take $\Sigma_j$ in $K(A;R,J)$.
For every $j\in[k]$, the triangles of $\Sigma_j$ form a connected triangulated surface of genus zero, namely a sphere with $1+|\{i:A_{ij}=1\}|$ holes, whose boundary is the disjoint union of the tether $\Lambda_j$ and the row cycles $C_i$ with $A_{ij}=1$.
Every edge of $K(A;R,J)$ that is not on a row cycle or a tether lies in either no triangle or exactly two triangles of $K(A;R,J)$, and in the latter case both triangles belong to the same set~$\Sigma_j$.
Any two triangles of $\Sigma_j$ are joined by a sequence of triangles of $\Sigma_j$ in which consecutive triangles share an edge that is not on a row cycle or a tether.
\end{lemma}

\begin{proof}
Each cone is a disk, and a sealed socket is the annulus $\mathrm{Cyl}(a,b)$.
A tubed socket between $a$ and $b$ with tube to $c$ is a pair of pants, which we check directly.
The punctured cylinder alone is not a surface, because it is pinched at~$b_0$; the tube repairs this.
In the tubed socket, the edges $a_1b_0$ and $b_0a_0$ of the hole lie in one triangle of $\mathrm{Cyl}^{\circ}(a,b)$ and in one triangle of the tube, the edge $a_0a_1$ lies only in the tube triangle $a_0a_1c_0$, and the \emph{link} of every vertex, the graph formed by the edges opposite to it in the triangles that contain it, is a single path; for instance, the link of $b_0$ is the path $b_1\,a_1\,c_1\,c_2\,a_0\,b_2$.
The edges in only one triangle form the three disjoint cycles $a$, $b$, and~$c$, and the socket has $9$ vertices, $21$ edges, and $11$ triangles, so its Euler characteristic is $-1$; hence it is a sphere with three holes.
The pieces of column~$j$ are glued in a chain along the complete triangle cycles $s^j_0,\dots,s^j_{d_j+1}$, so $\Sigma_j$ is a connected surface: a disk, followed by pairs of pants and annuli, followed by a disk.
Its boundary consists of the far ends of its tubes, namely $\Lambda_j$ and the row cycles $C_i$ with $A_{ij}=1$.
Each cone and each socket is connected through its interior edges, by direct inspection of the triangle lists, and consecutive pieces share the three edges of a cycle $s^j_t$, which are interior edges of~$\Sigma_j$; this gives the last claim.
All vertices of column $j$ other than those of the row cycles are private to column~$j$.
Hence an edge that lies on no row cycle and no tether has at least one private endpoint, so all triangles that contain it belong to one column $j$, and it is an interior edge of the surface $\Sigma_j$ or an edge in no triangle of~$K$.
An interior edge of a triangulated surface lies in exactly two of its triangles.
The fillings $D_i$ and $E_j$ use only edges of row cycles and tethers.
\end{proof}

\subsection{The crux: two-cycles are column dependencies}\label{sec:crux}

Over $\Ftwo$ we may ignore orientations.
For other fields we orient the triangles explicitly.
For a triangle cycle $x$ let $[x]=x_0x_1+x_1x_2+x_2x_0$ denote the $1$-cycle that traverses it in index order.
We orient the triangles of $\mathrm{Cyl}(a,b)$ as $a_t\,a_{t+1}\,b_t$ and $a_{t+1}\,b_{t+1}\,b_t$, for $t\in\Z_3$.
A direct check shows that at every edge shared by two of these triangles the two induced orientations are opposite, and that ${\bd\,\mathrm{Cyl}(a,b)=[a]-[b]}$.
The hole $h=a_0a_1b_0$ is one of these oriented triangles, so $\bd\,\mathrm{Cyl}^{\circ}(a,b)=[a]-[b]-[h]$, and the same rule gives $\bd\,\mathrm{Cyl}(h,c)=[h]-[c]$ for a tube.
Writing $o$ for the apex and $a$ for the coned cycle, we orient the cone on $a=s^j_0$ as $o\,a_{t+1}\,a_t$, so that its boundary is $-[s^j_0]$, and the cone on the last cycle $a=s^j_{d_j+1}$ as $o\,a_t\,a_{t+1}$, so that its boundary is $+[s^j_{d_j+1}]$.
Summing over the pieces of column~$j$, the contributions of the cycles $s^j_t$ and of the holes cancel, and the sum $[\Sigma_j]$ of the oriented triangles of $\Sigma_j$ satisfies $\bd[\Sigma_j]=-[\Lambda_j]-\sum_{i:\,A_{ij}=1}[C_i]$.
Since the pieces induce opposite orientations on every cycle along which they are glued, this orientation is coherent: the two triangles at any interior edge of $\Sigma_j$ induce opposite orientations on it.
We replace $[\Sigma_j]$ by its negative, and we orient the fillings so that $\bd D_i=[C_i]$ and $\bd E_j=[\Lambda_j]$.
Then
\begin{equation}\label{eq:surface-boundary}
  \bd[\Sigma_j]=[\Lambda_j]+\sum_{i:\,A_{ij}=1}[C_i].
\end{equation}
All signs in~\eqref{eq:surface-boundary} are equal; this is what makes the construction work over every field.

\begin{lemma}[Crux]\label{lem:crux}
Let $\F$ be a field, let $A$ be a $0/1$ matrix with $\supp(A)\subseteq P$, and let $R\subseteq[r]$ and $J\subseteq[k]$.
Then $K=K(A;R,J)$ satisfies
\[
  \beta_2(K;\F)\;=\;\dim Z_2(K;\F)\;=\;|J|-\rk_{\F}A[R,J].
\]
\end{lemma}

\begin{proof}
The first equality holds by~\eqref{eq:euler}, since $\dim K\le2$.
Let $z$ be a $2$-cycle of~$K$.
Consider an edge $e$ that lies in two triangles $\tau,\tau'$ of some~$\Sigma_j$ (Lemma~\ref{lem:surface}).
No other triangle contains~$e$, and $\tau,\tau'$ induce opposite orientations on~$e$, so the coefficient of $e$ in $\bd z$ is the difference of the coefficients of $\tau$ and $\tau'$ in~$z$.
Since $\bd z=0$, these coefficients agree.
By the last claim of Lemma~\ref{lem:surface}, $z$ has the same coefficient $x_j$ on every triangle of~$\Sigma_j$.
Every other triangle of $K$ is one of the fillings $D_i$ ($i\notin R$) or $E_j$ ($j\in J$).
Therefore
\[
  z=\sum_{j\in[k]}x_j[\Sigma_j]+\sum_{i\notin R}y_iD_i+\sum_{j\in J}w_jE_j
\]
for some scalars $x_j,y_i,w_j\in\F$.
By~\eqref{eq:surface-boundary},
\begin{align*}
  \bd z={}&\sum_{j\in J}(x_j+w_j)\,[\Lambda_j]+\sum_{j\notin J}x_j\,[\Lambda_j]\\
  &+\sum_{i\in R}\Big(\sum_{j}A_{ij}x_j\Big)[C_i]+\sum_{i\notin R}\Big(\sum_{j}A_{ij}x_j+y_i\Big)[C_i].
\end{align*}
The row cycles and the tethers have pairwise disjoint edge sets, so $\bd z=0$ holds if and only if all coefficients vanish.
For $j\notin J$ this forces $x_j=0$; for $j\in J$ it determines $w_j=-x_j$.
For $i\notin R$ it determines $y_i$, and for $i\in R$ it says $\sum_{j\in J}A_{ij}x_j=0$.
So $z\mapsto(x_j)_{j\in J}$ is a linear bijection from $Z_2(K)$ onto the kernel of $A[R,J]$, whose dimension is $|J|-\rk_{\F}A[R,J]$.
\end{proof}

Because $K(A;R,J)$ always contains the whole $1$-skeleton of $\cU(P)$, the number $\beta_0(K)$ and the numbers of vertices and edges do not change.
By~\eqref{eq:euler},
\begin{equation}\label{eq:b1-from-b2}
  \beta_1(K)=\beta_1(\cU^{(1)})+\beta_2(K)-n_2(K),
\end{equation}
where $\cU^{(1)}$ is the $1$-skeleton of~$\cU(P)$.
Here $\beta_1(\cU^{(1)})=\beta_0(\cU^{(1)})-n_0(\cU)+n_1(\cU)$ is the number of independent cycles of the graph $\cU^{(1)}$, a constant of the universe; intuitively, each triangle of $K$ either fills one of these cycles or creates a void.
So $\beta_1$ and $\beta_2$ determine each other once the number $n_2(K)$ of triangles is known.

\begin{lemma}[Locality]\label{lem:locality}
Let $K=K(A;R,J)$.
\begin{enumerate}[label=(\roman*)]
\item If $A'$ is a $0/1$ matrix that differs from $A$ in exactly one entry $(i,j)\in P$, then $K(A';R,J)$ differs from $K$ in exactly $7$ triangles.
\item If $R'=R\triangle\{i\}$ and $J'=J\triangle\{j\}$, then $K(A;R',J)$ and $K(A;R,J')$ each differ from $K$ in exactly one triangle.
\end{enumerate}
In both cases the change can be carried out by single-triangle insertions and deletions in any order, and every intermediate family of simplices is a subcomplex of~$\cU(P)$.
\end{lemma}

\begin{proof}
Flipping $A_{ij}$ exchanges the seal of the socket of position $(i,j)$ (one triangle) for its tube (six triangles), or vice versa: one triangle is deleted and six are inserted, or the reverse.
Switching row $i$ or column $j$ inserts or deletes the filling $D_i$ or~$E_j$.
All edges and vertices of $\cU(P)$ are always present, so every set of triangles together with them is a subcomplex.
\end{proof}

\subsection{The equivalence}\label{sec:equiv-thm}

\begin{theorem}[Rank-completeness]\label{thm:equiv}
Let $\F$ be a field.
\begin{enumerate}[label=(\roman*)]
\item For every fixed $p\ge0$, dynamic $\beta_p$ over $\F$ reduces to dynamic rank under switches over $\F$ of a $\{0,\pm1\}$-matrix with overhead $(1,1)$, and to dynamic rank over $\F$ of $\{0,\pm1\}$-matrices with overhead $(p+2,1)$.
\item\label{it:equiv-ii} Dynamic rank over $\F$ of $0/1$ matrices reduces to dynamic $\beta_2$ over $\F$ of $2$-complexes, and to dynamic $\beta_1$ over $\F$ of $2$-complexes, with overhead $(7,1)$. Every entry flip becomes $7$ simplex updates and every switch becomes one simplex update.
\item For every fixed $p\ge1$, dynamic $\beta_1$ over $\F$ of $2$-complexes reduces to dynamic $\beta_{p}$ over $\F$ of $(p+1)$-complexes with overhead $(3^{p-1},1)$, and for every fixed $p\ge2$, dynamic $\beta_2$ over $\F$ of $2$-complexes reduces to dynamic $\beta_p$ over $\F$ of $p$-complexes with overhead $(3^{p-2},1)$.
\item Dynamic rank over $\F$ of $\{0,\pm1\}$-matrices reduces to dynamic $\beta_2$ over $\F$ of $2$-complexes, and to dynamic $\beta_1$ over $\F$ of $2$-complexes, with overhead $(8,1)$.
\end{enumerate}
\end{theorem}

Since boundary matrices have entries in $\{0,\pm1\}$, Theorem~\ref{thm:equiv} has the following consequence.

\begin{corollary}\label{cor:equiv}
Let $\F$ be a field, and let $p\ge1$ and $d\ge\max\{p,2\}$ be fixed.
Dynamic rank over $\F$ of $\{0,\pm1\}$-matrices, dynamic $\beta_1$ over $\F$ of $2$-complexes, and dynamic $\beta_p$ over $\F$ of complexes of dimension at most~$d$ reduce to one another with constant overhead.
Over $\Ftwo$ and over $\mathbb{F}_3$, this includes dynamic rank of all matrices.
\end{corollary}

Part~(iii) uses the suspension.
The \emph{suspension} of a complex $K$ with two new vertices $a,b$ is
\[
  SK=K\cup\{\{a\},\{b\}\}\cup\{\sigma\cup\{a\},\,\sigma\cup\{b\}:\sigma\in K\}.
\]

\begin{lemma}\label{lem:suspension}
For every nonempty subcomplex $K$ of a universe $\cU$, every field~$\F$, and every $p\ge0$, we have $\widetilde H_{p+1}(SK;\F)\cong\widetilde H_p(K;\F)$; moreover, $\beta_{p+1}(SK)=\beta_p(K)$ for every subcomplex $K$, empty or not, and every $p\ge1$.
Inserting or deleting a simplex $\sigma$ of $K$ corresponds to inserting or deleting the three simplices $\sigma$, $\sigma\cup\{a\}$, $\sigma\cup\{b\}$ of $SK$, in an order in which every intermediate family is a subcomplex of~$S\cU$.
\end{lemma}

\begin{proof}
If $K=\emptyset$, then $SK$ consists of two points, and $\beta_{p+1}(SK)=0=\beta_p(K)$ for $p\ge1$.
Otherwise $SK$ is the union of the cones $aK=\{\{a\}\}\cup K\cup\{\sigma\cup\{a\}:\sigma\in K\}$ and $bK$, which are contractible, and $aK\cap bK=K$.
The Mayer--Vietoris sequence in reduced homology~\cite[Section~2.2]{Hat02} therefore gives $\widetilde H_{p+1}(SK)\cong\widetilde H_p(K)$ for every $p\ge0$, and reduced and unreduced homology agree in positive degrees.
For insertion, insert $\sigma$ first and then its two cones; for deletion, use the reverse order. Every intermediate family is then a subcomplex of~$S\cU$.
\end{proof}

\begin{proof}[Proof of Theorem~\ref{thm:equiv}]
Part~(i) is Lemma~\ref{lem:betti-to-rank}.
The rank instance is the block matrix of Lemma~\ref{lem:betti-to-rank}; its size is $O(n)$ for fixed~$p$.
In the version with switches, the instance is the block matrix itself, preprocessing switches off all its columns, and afterwards only switches are used.
In the version with entry updates, the admissible positions are the at most ${(p+1)\,n_p(\cU)}+{(p+2)\,n_{p+1}(\cU)}$ nonzero positions of the block matrix, and all entries are initially zero.
Inserting or deleting a $q$-simplex for $q\in\{p,p+1\}$ switches one column, or updates its at most $q+1$ nonzero entries; the query subtracts the answer from a counter for~$n_p(K)$.

For part~(ii), we map an instance $(r,k,P)$ to the universe $\cU(P)$ of Definition~\ref{def:gadget}.
In preprocessing we insert the simplices of $K(0;[r],[k])$ in order of increasing dimension.
We maintain the invariant $K=K(A;R,J)$, together with counters for $|J|$ and $n_2(K)$.
By Lemma~\ref{lem:locality}, a flip is carried out by $7$ simplex updates and a switch by one.
By Lemma~\ref{lem:crux}, $\rk A[R,J]=|J|-\beta_2(K)$, and by~\eqref{eq:b1-from-b2}, $\beta_2(K)=\beta_1(K)-\beta_1(\cU^{(1)})+n_2(K)$, so one query of either Betti number answers a rank query.
The universe has $O(r+k+|P|)$ simplices.

For part~(iii) the universe is the $(p-1)$-fold suspension $S^{p-1}\cU$ of~$\cU$ (the suspension applied $p-1$ times, with two new vertices each time), which has fewer than $3^{p-1}(n+1)$ simplices.
In preprocessing we insert the $3^{p-1}-1$ simplices of the $(p-1)$-fold suspension of the empty complex, in order of increasing dimension.
Afterwards each update of $K$ becomes $3^{p-1}$ updates, by Lemma~\ref{lem:suspension} applied $p-1$ times.
The reduction from $\beta_2$ of $2$-complexes to $\beta_p$ of $p$-complexes is the same with the $(p-2)$-fold suspension.

For part~(iv), let $\cU^{\pm}(P)$ consist of $\cU(P)$ and, for every $(i,j)\in P$, the \emph{reversed tube} $\mathrm{Cyl}(h,\bar C_i)$ with its edges, where $h$ is the hole of the socket of $(i,j)$ and $\bar C_i=(c_0,c_2,c_1)$ traverses $C_i=(c_0,c_1,c_2)$ backwards.
For $A\in\{0,\pm1\}^{r\times k}$ with $\supp(A)\subseteq P$, the complex $K^{\pm}(A;R,J)$ is defined as in Definition~\ref{def:gadget}, with all vertices and edges of $\cU^{\pm}(P)$, and with the reversed tube instead of the tube at every position with $A_{ij}=-1$.
A socket with a reversed tube is a socket tubed to the triangle cycle $\bar C_i$, so Lemma~\ref{lem:surface} holds with $A_{ij}=1$ replaced by $A_{ij}\neq0$.
Betti numbers do not depend on orientations, so we may choose them separately for each complex: every triangle of $K^{\pm}(A;R,J)$ is oriented as the piece of $K^{\pm}(A;R,J)$ that contains it orients it.
Since $\bd\,\mathrm{Cyl}(h,\bar C_i)=[h]+[C_i]$, the computation before~\eqref{eq:surface-boundary} gives $\bd[\Sigma_j]=[\Lambda_j]+\sum_iA_{ij}[C_i]$, and the proof of Lemma~\ref{lem:crux} gives $\beta_2(K^{\pm}(A;R,J);\F)=|J|-\rk_{\F}A[R,J]$.
Changing an entry between $0$ and $\pm1$ exchanges the seal and a tube, which takes $7$ triangle updates; changing it between $1$ and $-1$ deletes $4$ triangles and inserts~$4$, because the two tubes share the two triangles $h_0h_1c_0$ and $h_2c_1c_2$.
The universe still has $O(r+k+|P|)$ simplices, and the rest of the proof is as for part~(ii).
\end{proof}

\begin{proof}[Proof of Corollary~\ref{cor:equiv}]
By Theorem~\ref{thm:equiv}(i), dynamic $\beta_p$ reduces to dynamic rank of the block matrix of Lemma~\ref{lem:betti-to-rank}, whose entries lie in $\{0,\pm1\}$.
By part~(iv), dynamic rank of $\{0,\pm1\}$-matrices reduces to dynamic $\beta_1$ and to dynamic $\beta_2$ of $2$-complexes.
For $p\le2$ we are done, since $d\ge2$; for $p\ge3$, part~(iii) reduces dynamic $\beta_2$ of $2$-complexes to dynamic $\beta_p$ of $p$-complexes, and $p\le d$.
Over $\Ftwo$ and over $\mathbb{F}_3$, every matrix has entries in $\{0,\pm1\}$.
\end{proof}

Two consequences are worth stating separately.
First, composing~(ii) and~(iv) with the switch-only version of~(i) for $p=2$ shows that switches alone are as powerful as entry updates.

\begin{corollary}\label{cor:switch-only}
Let $\F$ be a field.
Dynamic rank over $\F$ of $0/1$ matrices reduces with overhead $(7,1)$, and dynamic rank over $\F$ of $\{0,\pm1\}$-matrices reduces with overhead $(8,1)$, to dynamic rank under switches over $\F$ of a $\{0,\pm1\}$-matrix with $O(N)$ nonzero entries, using column switches only.
Over $\Ftwo$ and over $\mathbb{F}_3$, this applies to all matrices.
\end{corollary}

Second, if nothing is known in advance about the positions of the nonzero entries, we may take $P=[r]\times[k]$; then the universe has $\Theta(rk)$ simplices.
This is the setting of the lower bounds in Section~\ref{sec:lower}.

